\begin{proposition}[6.10.5]
\label{IV.2.s6.10.5-fr}
Con las notaciones de \textup{(6.10.1)}, supongamos que $Y$ es conexo y no
vacío, y que $\mathcal F$ es normalmente plano a lo largo de $Y$. Para todo
entero $n\geq0$, se pone
\begin{equation}
r(n)=\operatorname{rg}_{\mathcal O_Y}
(\operatorname{gr}_{\mathcal J}^{n}(\mathcal F))
\tag{6.10.5.1}\label{IV.2.s6.10.5.1-fr}
\end{equation}
(el $\mathcal O_Y$-Módulo localmente libre
$\operatorname{gr}_{\mathcal J}^{n}(\mathcal F)$ tiene necesariamente rango
constante). Entonces:
\begin{enumerate}
\item[\textup{(i)}] Existe un polinomio $P\in\mathbf Q[T]$ tal que $P(n)=r(n)$ para
todo $n$ suficientemente grande.
\item[\textup{(ii)}] Se tiene $Y\cap\operatorname{Supp}(\mathcal F)=\varnothing$
(es decir, $\mathcal F/\mathcal J\mathcal F=0$), o bien
$Y\cap\operatorname{Supp}(\mathcal F)=Y$ (es decir,
$Y\subset\operatorname{Supp}(\mathcal F)$).
\end{enumerate}
En el segundo caso, denotemos por $d\setmin 1$ el grado de $P$; para todo punto maximal $y$
de $Y$ se tiene
\begin{equation}
\dim(\mathcal F_y)=d
\tag{6.10.5.2}\label{IV.2.s6.10.5.2-fr}
\end{equation}
y, en particular,
\begin{equation}
\operatorname{codim}(Y,\operatorname{Supp}(\mathcal F))=d.
\tag{6.10.5.3}\label{IV.2.s6.10.5.3-fr}
\end{equation}
\begin{enumerate}
\item[\textup{(iii)}] Supongamos que $Y\subset\operatorname{Supp}(\mathcal F)$. Para todo
$x\in Y$ se tiene
\end{enumerate}
\begin{equation}
\dim(\mathcal F_x)=\dim(\mathcal O_{Y,x})+d.
\tag{6.10.5.4}\label{IV.2.s6.10.5.4-fr}
\end{equation}

Más precisamente, existen en $\mathcal J_x$ elementos $f_i$
$(1\leq i\leq d)$ que forman parte de un sistema de parámetros para
$\mathcal F_x$ \textup{(\textbf{0}, 16.3.6)} y tales que, en
$\operatorname{Spec}(\mathcal O_{X,x})$, se tiene
$V(\mathcal J_x)=V(\sum_{i=1}^{d}f_i\mathcal O_{X,x})\cap\operatorname{Supp}(\mathcal F)$.
\end{proposition}

\oldpage[IV-2]{157}
\begin{proof}
Como se supone que $Y$ es conexo, la primera afirmación de (ii) se sigue de
(6.10.3, (i)). Si se tiene
$Y\cap\operatorname{Supp}(\mathcal F)=\varnothing$, la afirmación (i) es
trivial con $P=0$. Supongamos que
$\operatorname{Supp}(\mathcal F)\supset Y$, y sea $y$ un punto maximal de
$Y$; como entonces $Y=\operatorname{Supp}(\mathcal F/\mathcal J\mathcal F)$,
$\mathcal F_y/\mathcal J_y\mathcal F_y$ es un
$\mathcal O_{X,y}$-módulo de longitud finita (3.1.4). Poniendo
$s(n)=\operatorname{long}(\operatorname{gr}_{\mathcal J_y}^{n}(\mathcal F_y))$,
se sabe que existe un polinomio $R\in\mathbf Q[T]$ tal que $s(n)=R(n)$ para
$n$ suficientemente grande, a saber, el polinomio $H(n+1)\setmin H(n)$, donde $H$ es el polinomio
de Hilbert-Samuel de $\mathcal F_y$ para la filtración
$\mathcal J_y$-preádica (\textbf 0, 16.2.1). Como los
$\mathcal O_{Y,y}$-módulos
$\operatorname{gr}_{\mathcal J_y}^{n}(\mathcal F_y)$ son libres por
hipótesis, se tiene $r(n)=s(n)/m$, donde $m$ designa la longitud del
$\mathcal O_{X,y}$-módulo $\mathcal O_{Y,y}=\mathcal O_{X,y}/\mathcal J_y$; se satisface,
pues, (i)
tomando $P=\frac{1}{m}R$. Se sabe, además, (\textbf 0, 16.2.3) que
$\deg(H)=\dim(\mathcal F_y)$, de donde se sigue la relación (6.10.5.2); la relación
(6.10.5.3) se deduce de ella mediante (5.1.12.2).

Queda por demostrar (iii) para un punto cualquiera $x\in Y$. Pongamos
$A=\mathcal O_{X,x}$, $\mathfrak J=\mathcal J_x$, de modo que
$\mathcal O_{Y,x}=A/\mathfrak J$, y $M=\mathcal F_x$; sea
$S=\operatorname{gr}_{\mathfrak J}^{\bullet}(A)$, que es una
$(A/\mathfrak J)$-álgebra graduada de tipo finito, en grados no negativos, tal
que $S_0=A/\mathfrak J$, y generada por sus elementos homogéneos de grado
$1$; sea, finalmente,
$N=\operatorname{gr}_{\mathfrak J}^{\bullet}(M)$, que es un $S$-módulo
graduado finitamente generado, cuya componente homogénea $N_n$ es, por hipótesis,
un $(A/\mathfrak J)$-módulo libre de longitud $r(n)$. Sean $\mathfrak m$
el ideal maximal de $A$, $k=A/\mathfrak m$ su cuerpo residual;
$B=S\otimes_A k$ es una $k$-álgebra graduada de tipo finito, en grados
no negativos, generada por sus elementos homogéneos de grado $1$ y tal que
$B_0=k$, de modo que $\mathfrak q=B^+=\bigoplus_{n\geq1}B_n$ es un ideal
maximal de $B$; $E=N\otimes_A k$ es un $B$-módulo graduado finitamente generado
tal que $\operatorname{rg}_k(E_n)=r(n)$. Apliquemos (\textbf 0, 16.2.7) al anillo
graduado $B=S\otimes_A k$ y al $B$-módulo graduado $E=N\otimes_A k$, y sea
$f_i\in\mathfrak J^{n_i}$ $(1\leq i\leq d)$ un elemento cuya imagen es $t_i$
en $B_{n_i}$. Para $n\geq\sup_i(n_i)$, consideremos el sub-$A$-módulo
$\sum_{i=1}^{d}f_i\mathfrak J^{n-n_i}M$ de $M$; como la componente homogénea
de grado $n$ del submódulo $\sum_{i=1}^{d}t_iE$ de $E$ es igual a
$E_n$ cuando $n$ es suficientemente grande, se ve que, para $n$ suficientemente grande, se tiene
\[
\mathfrak J^nM=\sum_{i=1}^{d}f_i\mathfrak J^{n-n_i}M+
\mathfrak m\mathfrak J^nM
\]
y como $\mathfrak J^nM$ es un $A$-módulo finitamente generado, se sigue,
por el lema de Nakayama, que
\[
\mathfrak J^nM=\sum_{i=1}^{d}f_i\mathfrak J^{n-n_i}M
\subset\sum_{i=1}^{d}f_iM.
\]
Si $\mathfrak a$ es el anulador de $M$, se tiene, pues, (\textbf{0}$_{\rm I}$, 1.7.5)
en $\operatorname{Spec}(A)$
\[
V\left(\sum_{i=1}^{d}f_iA\right)\cap V(\mathfrak a)
\subset V(\mathfrak J^n)\cap V(\mathfrak a)
=V(\mathfrak J)\cap V(\mathfrak a)=V(\mathfrak J)
\]
ya que, por hipótesis,
$Y\cap\operatorname{Spec}(A)=V(\mathfrak J)\cap\operatorname{Supp}(M)
=V(\mathfrak J)$. Como, por otra parte, los $f_i$ pertenecen a
$S$, se tiene
$V(\mathfrak J)\subset V(\sum_{i=1}^{d}f_iA)$, lo que demuestra la última
relación de (iii). Queda por probar que los $f_i$ forman parte de un sistema
de parámetros para $M$. Reemplazando $A$ por $A/\mathfrak a$ y los $f_i$ por
sus imágenes en $A/\mathfrak a$, es posible limitarse al caso en que
$\mathfrak a=0$; se acaba de ver que se tiene
$\dim(A/\mathfrak J)=\dim(A/(\sum_{i=1}^{d}f_iA))$
y solo queda, pues, demostrar (\textbf 0, 16.3.7) que se tiene
\[
\dim(A)\geq\dim(A/\mathfrak J)+d.
\]

Ahora bien, sea $\mathfrak p$ un ideal primo de $A$ que contiene a $\mathfrak J$ y
tal que $\dim(A/\mathfrak J)=\dim(A/\mathfrak p)$, de modo que
$\mathfrak p$ es minimal entre los ideales primos que contienen a
$\mathfrak J$. Se tiene, por tanto, $\mathfrak p=\mathfrak j_y$, donde
$y\in\operatorname{Spec}(A)$ es un punto maximal de $Y$. Pero, en virtud de
(6.10.5.2) y de la hipótesis
$\operatorname{Spec}(A)=\operatorname{Supp}(M)$, se tiene
$\dim(A_{\mathfrak p})=d$; la desigualdad
$\dim(A/\mathfrak p)+\dim(A_{\mathfrak p})\leq\dim(A)$ (\textbf 0, 16.1.4)
concluye la demostración.
\end{proof}
